\documentclass[10pt,a4paper]{amsart}
\usepackage[margin=19mm]{geometry}
\usepackage{amsmath,amssymb,mathtools}
\usepackage[scr=rsfs,cal=euler]{mathalfa}
\usepackage{xfrac}
\usepackage[unicode,hidelinks,bookmarks=false]{hyperref}
\newcommand{\ie}{i.e.\ }
\newcommand{\cf}{cf.\ }
\newcommand{\resp}{respectively\ }
\newcommand{\Subsection}{\subsection}
\providecommand{\nobreakdash}{\nobreak\hbox{-}}
\setlength{\emergencystretch}{2em}
\multlinegap0pt
\usepackage{bm}
\usepackage{bbm}
\usepackage{dsfont}
\usepackage{xspace}
\usepackage{cancel}
\usepackage{mathtools}
\usepackage{amsthm}
\makeatletter
\@ifundefined{theorem}{\newtheorem{theorem}{Theorem}}{}
\makeatother
\makeatletter
\expandafter\ifx\csname Example\endcsname\relax
\newenvironment{Example}{\begin{example}\rm}{\end{example}}
\fi
\expandafter\ifx\csname Examples\endcsname\relax
\newenvironment{Examples}{\begin{examples}\rm}{\end{examples}}
\fi
\newcommand{\Ho}{\mathcal H}
\expandafter\ifx\csname Remark\endcsname\relax
\newenvironment{Remark}{\begin{remark}\rm}{\end{remark}}
\fi
\expandafter\ifx\csname Remarks\endcsname\relax
\newenvironment{Remarks}{\begin{foo}\rm}{\end{foo}}
\fi
\newcommand{\Tor}{\mathrm{Tor}}
\newcommand{\V}{\mathcal V}
\makeatother
\title[Sub-Laplacian comparison theorems on Riemannian foliations w]{Sub-Laplacian comparison theorems on Riemannian foliations with minimal leaves and applications}
\author{Fabrice Baudoin}
\date{Source: arXiv:2509.13276v1 (CC BY 4.0)}
\begin{document}
\maketitle

\noindent\emph{Source and licence.} Sub-Laplacian comparison theorems on Riemannian foliations with minimal leaves and applications. Version arXiv:2509.13276v1 (CC BY 4.0), distributed under the Creative Commons Attribution 4.0 International licence, as recorded in the arXiv licence metadata for this exact version and shown on the abstract page https://arxiv.org/abs/2509.13276v1. Full rights notice: licenses/arxiv-2509.13276-CC-BY-4.0.txt.\par\smallskip

\noindent\textit{Editorial scope.} Counted unit: the Theorem whose source label is \texttt{GB total} in the article (source0.tex lines 1182--1191, source label \texttt{GB total}), delivered together with the article's own complete original proof of it (source0.tex lines 1193--1265, one \texttt{proof} environment of 73 lines). No mathematics is written, repaired or re-derived: statement and proof are the article's own text line for line. Required context retained uncounted: geodesic1 (lines 470--472); comparison laplacian (lines 805--822). Every definition, display and label of those lines is retained unchanged. Omitted from this package and not redistributed: 1--469, 473--804, 823--1181, 1192--1192, 1266--1435. Nothing inside a retained excerpt is omitted or altered: every retained body is delivered exactly as the article's lines are written, so no omission receipt of any kind accompanies this package. Citations: the retained excerpts carry no citation command at all, so no bibliography entry of the article is transcribed and no reference is redistributed. Reference adaptation: none was needed and none was made; every reference inside the delivered unit resolves inside it, so no unresolved reference or citation is produced. Printed slips kept literal: no printed word, symbol, hypothesis or number of the counted statement or of its proof was altered. Layout and class: the article's own class is \texttt{article} with packages including amsfonts, amsmath, amsthm, amssymb, enumerate, color, cancel, authblk, enumitem; the wrapper is the lane's pinned \texttt{amsart} editorial wrapper at 10pt on A4, which restates the theorem-like environments the retained text needs and adds the article's own macro definitions that the retained text needs (\texttt{\textbackslash Ho}, \texttt{\textbackslash Tor}, \texttt{\textbackslash V}), with extra wrapper packages (bm, bbm, dsfont, xspace, cancel, mathtools). The delivered numbering is the unit's own. Assets: no retained span contains a figure environment, a graphics-inclusion command, a PDF-page inclusion, a table or a TikZ picture; no image file is copied into this package, the package declares no asset dependency and no image file is redistributed with it. Licence: version arXiv:2509.13276v1 of this article is distributed under the Creative Commons Attribution 4.0 International licence, as recorded in the arXiv licence metadata for this exact version and shown on the abstract page https://arxiv.org/abs/2509.13276v1. A full rights notice is delivered at licenses/arxiv-2509.13276-CC-BY-4.0.txt. Characters: every retained excerpt is pure ASCII, so the delivered engine, XeLaTeX with its default Latin Modern text font, reports no missing character.
\begin{align}\label{geodesic1}
\nabla_{\dot\gamma} \dot\gamma+ J_{\dot\gamma} \dot\gamma =0.
\end{align}

\begin{theorem}\label{comparison laplacian}
Let $p \in M$ and denote $r_p(x)=d(p,x)$. Assume that there exists $K \in \mathbb{R}$ such that for every $U \in \mathfrak{X}(TM)$
\[
\mathfrak{R}(U,U) \ge K |U|^2.
\] 
Then, for $x\neq p$ not in the  cut-locus of $p$,
\begin{align*}
\Delta_\Ho r_p (x) \le 
\begin{cases}
 \sqrt { nK} \cot\left(\sqrt { \frac{K}{n}} r_p(x) \right) &\text{if}\ K>0,
\\
\displaystyle\frac{n}{r_p(x)} &\text{if}\ K = 0,
\\
 \sqrt{n |K|} \coth\left(\sqrt{\frac{|K|}{n}} r_p(x)\right) &\text{if}\ K<0.
\end{cases}
\end{align*}

\end{theorem}

\begin{theorem}\label{GB total}
 Assume that the foliation is totally geodesic and that there exists $K \in \mathbb{R}$ such that for every $U \in \mathfrak{X}(TM)$
\[
\mathfrak{R}(U,U) -\frac{1}{4} (J_U,J_U)_\Ho \ge K |U|^2.
\] 
Then for every $f \in \mathbf{Lip}_b (M) \cap C^2(M)$ and $t \ge 0$, $P_tf \in \mathbf{Lip}_b (M)$ and we have
\[
|\nabla P_t f|(x)\le e^{-Kt} P_t |\nabla f| (x), \quad x \in M.
\] 
\end{theorem}

\begin{proof}
We use a  different coupling. The idea is to consider the following connection
\[
\overset{\circ}{\nabla}_X Y=\nabla_XY +J_XY.
\]
Since the foliation is totally geodesic one can check that $\overset{\circ}{\nabla}$ satisfies the following properties which will be key in the following arguments:
\begin{enumerate}
\item $\overset{\circ}{\nabla}$ is a metric connection;
\item For every $X \in \mathfrak{X}(\Ho)$ and $Y \in \mathfrak{X}(TM)$, $\overset{\circ}{\nabla}_Y X \in \mathfrak{X}(\Ho)$;
\item The geodesic equation can be written $\overset{\circ}{\nabla}_{\dot \gamma} \dot\gamma=0$, see \eqref{geodesic1}.
\end{enumerate}
For $p,q \in \mathcal{C}$ define the transport map $\overset{\circ}{\mathfrak{P}}_{p \to q}:T_pM \to T_qM$ given for $v \in T_pM$ by $\overset{\circ}{\mathfrak{P}}_{p \to q}(v)=\overset{\circ}{v}$ where $\overset{\circ}{v}$ is the $\overset{\circ}{\nabla}$-parallel transport of $v$ along the unit speed geodesic $\gamma:[0,r] \to M$ with endpoints $p,q$. As before, consider  a horizontal Brownian motion $\xi_t$ started from $p \in M$. For $q \in M$ such that $(p,q) \in \mathcal{C}$ we construct the process $\tilde{\xi}_t$ that solves the stochastic differential equation
\begin{align*}
\begin{cases}
d\tilde{\xi}_t=\overset{\circ}{\mathfrak{P}}_{p \to q} \circ d\xi_t \\
\tilde{\xi}_0=q.
\end{cases}
\end{align*}
As before, we will  ignore the genuine but inessential difficulties related to cut-locus and assume that the couple $(\xi_t,\tilde{\xi}_t)$ is defined for all $t$'s. 

By construction, the process is a diffusion on $M \times M$ with generator
\begin{align*}
\overset{\circ}{\Delta}\,^c_{\Ho,\Ho} f (p,q)=\Delta^1_{\Ho} f (p,q)+\Delta^2_{\Ho} f (p,q) +2 \sum_{i=1}^n v_i^1 \overset{\circ}{v}\,^2_i f(p,q).
\end{align*}
In It\^o's formula the local martingale part of $d(\xi_t,\tilde{\xi}_t)$ is the same as the local martingale part of
\begin{align}\label{quadratic part }
\int_0^t \left\langle \nabla^1 d(\xi_t,\tilde{\xi}_t), \circ d\xi_t \right\rangle +\left\langle \nabla^2 d(\xi_t,\tilde{\xi}_t), \circ d\tilde{\xi}_t \right\rangle
\end{align}
which vanishes since $\nabla^1 d(\xi_t,\tilde{\xi}_t) +\overset{\circ}{\mathfrak{P}}_{\tilde{\xi}_t \to \xi_t}\nabla^2 d(\xi_t,\tilde{\xi}_t)=0$ from the fact that any geodesic satisfies $\overset{\circ}{\nabla}_{\dot \gamma} \dot\gamma=0$.


We now need an estimate on $\overset{\circ}{\Delta}\,^c_{\Ho,\Ho} d (p,q)$. The argument is essentially the same as in the proof of Theorem \ref{comparison laplacian}. From the index lemma one has
\[
\overset{\circ}{\Delta}\,^c_{\Ho,\Ho} d (p,q) \le  \sum_{i=1}^n I(\gamma,Y_i)
\]
where the $Y_i$'s are arbitrary vector fields along $\gamma$ such that
\[
Y_i(0)=v_i, \, \, Y_i(r)=\overset{\circ}{v}_i.
\]
Pick
\[
Y_i(t)=\mathfrak{c}_K (t) X_i(t),
\]
where $X_i$ is the $\overset{\circ}{\nabla}$-parallel transport of $v_i$ along $\gamma$ and $\mathfrak{c}_K$ is as in the proof of Theorem \ref{comparison laplacian}. We then compute
\begin{align*}
 \sum_{i=1}^n I(\gamma,Y_i)  &=  \sum_{i=1}^n  \int_0^r  \left | \nabla_{\dot \gamma}  Y_i +\frac{1}{2} J_{\dot \gamma}  Y_i   \right|^2  -\left \langle \mathrm{Riem}^\nabla(\dot \gamma,Y_i)Y_i, \dot \gamma \right \rangle  +\left\langle (\nabla_{Y_i} \Tor^\nabla)(Y_i,\dot\gamma),\dot \gamma_\V \right\rangle\\
 &\quad -\frac{1}{4} | J_{\dot \gamma} Y_i |^2_\Ho+ |\Tor^\nabla(Y_i,\dot\gamma)|^2 \, dt \\
 &=\sum_{i=1}^n  \int_0^r  \mathfrak{c}'_K(t)^2  - \mathfrak{c}_K(t)^2 \left(  \left \langle \mathrm{Riem}^\nabla(\dot \gamma,X_i)X_i, \dot \gamma \right \rangle  +\left\langle (\nabla_{X_i} \Tor^\nabla)(X_i,\dot\gamma),\dot \gamma_\V \right\rangle \right. \\
 &\left. \quad + |\Tor^\nabla(X_i,\dot\gamma)|^2 \right) \, dt \\
 &=  \int_0^r n  \mathfrak{c}'_K(t)^2-\mathfrak{c}_K(t)^2\left( \mathfrak{R}(\dot \gamma,\dot \gamma) +\frac{1}{4} (J_{\dot \gamma},J_{\dot \gamma})_\Ho \right)dt \\
 &\le  \int_0^r n  \mathfrak{c}'_K(t)^2- K \mathfrak{c}_K(t)^2  dt \\
 &\le -Kr.
\end{align*}
Therefore, we have
\[
d(\xi_t,\tilde{\xi}_t) \le d(p,q)-K \int_0^t d(\xi_s,\tilde{\xi}_s) ds
\]
which yields
\[
d(\xi_t,\tilde{\xi}_t) \le  d(p,q) e^{-Kt}.
\]
Let now $f \in \mathbf{Lip}_b (M)\cap C^2(M)$. One has
\begin{align*}
|P_tf(q)-P_tf(p)|&= \left| \mathbb{E} \left( f(\tilde{\xi}_t) \right)- \mathbb{E}\left(f(\xi_t)\right) \right| \\
 &\le \mathbb{E} \left( \left| f(\tilde{\xi}_t)- f(\xi_t)\right| \right) \\
 &\le   \mathbb{E} \left(d(\xi_t,\tilde{\xi}_t)  \sup_{\rho \in B(\xi_t,d(p,q) e^{-Kt})} \left| \nabla f (\rho)\right| \right) \\
 &\le d(p,q) e^{-Kt}\mathbb{E} \left(  \sup_{\rho \in B(\xi_t,d(p,q) e^{-Kt})} \left| \nabla f (\rho)\right| \right).
\end{align*} 
Using the density of $\mathcal{C}$ in $M\times M$ and letting $q \to p$ yields
\[
|\nabla P_t f|(p)\le e^{-Kt} P_t |\nabla f| (p).
\]
\end{proof}

\end{document}
